% Options for packages loaded elsewhere
% Options for packages loaded elsewhere
\PassOptionsToPackage{unicode}{hyperref}
\PassOptionsToPackage{hyphens}{url}
\PassOptionsToPackage{dvipsnames,svgnames,x11names}{xcolor}
%
\documentclass[
  spanish,
  letterpaper,
  DIV=11,
  numbers=noendperiod]{scrartcl}
\usepackage{xcolor}
\usepackage[margin=0cm]{geometry}
\usepackage{amsmath,amssymb}
\setcounter{secnumdepth}{5}
\usepackage{iftex}
\ifPDFTeX
  \usepackage[T1]{fontenc}
  \usepackage[utf8]{inputenc}
  \usepackage{textcomp} % provide euro and other symbols
\else % if luatex or xetex
  \usepackage{unicode-math} % this also loads fontspec
  \defaultfontfeatures{Scale=MatchLowercase}
  \defaultfontfeatures[\rmfamily]{Ligatures=TeX,Scale=1}
\fi
\usepackage{lmodern}
\ifPDFTeX\else
  % xetex/luatex font selection
\fi
% Use upquote if available, for straight quotes in verbatim environments
\IfFileExists{upquote.sty}{\usepackage{upquote}}{}
\IfFileExists{microtype.sty}{% use microtype if available
  \usepackage[]{microtype}
  \UseMicrotypeSet[protrusion]{basicmath} % disable protrusion for tt fonts
}{}
\makeatletter
\@ifundefined{KOMAClassName}{% if non-KOMA class
  \IfFileExists{parskip.sty}{%
    \usepackage{parskip}
  }{% else
    \setlength{\parindent}{0pt}
    \setlength{\parskip}{6pt plus 2pt minus 1pt}}
}{% if KOMA class
  \KOMAoptions{parskip=half}}
\makeatother
% Make \paragraph and \subparagraph free-standing
\makeatletter
\ifx\paragraph\undefined\else
  \let\oldparagraph\paragraph
  \renewcommand{\paragraph}{
    \@ifstar
      \xxxParagraphStar
      \xxxParagraphNoStar
  }
  \newcommand{\xxxParagraphStar}[1]{\oldparagraph*{#1}\mbox{}}
  \newcommand{\xxxParagraphNoStar}[1]{\oldparagraph{#1}\mbox{}}
\fi
\ifx\subparagraph\undefined\else
  \let\oldsubparagraph\subparagraph
  \renewcommand{\subparagraph}{
    \@ifstar
      \xxxSubParagraphStar
      \xxxSubParagraphNoStar
  }
  \newcommand{\xxxSubParagraphStar}[1]{\oldsubparagraph*{#1}\mbox{}}
  \newcommand{\xxxSubParagraphNoStar}[1]{\oldsubparagraph{#1}\mbox{}}
\fi
\makeatother


\usepackage{longtable,booktabs,array}
\newcounter{none} % for unnumbered tables
\usepackage{calc} % for calculating minipage widths
% Correct order of tables after \paragraph or \subparagraph
\usepackage{etoolbox}
\makeatletter
\patchcmd\longtable{\par}{\if@noskipsec\mbox{}\fi\par}{}{}
\makeatother
% Allow footnotes in longtable head/foot
\IfFileExists{footnotehyper.sty}{\usepackage{footnotehyper}}{\usepackage{footnote}}
\makesavenoteenv{longtable}
\usepackage{graphicx}
\makeatletter
\newsavebox\pandoc@box
\newcommand*\pandocbounded[1]{% scales image to fit in text height/width
  \sbox\pandoc@box{#1}%
  \Gscale@div\@tempa{\textheight}{\dimexpr\ht\pandoc@box+\dp\pandoc@box\relax}%
  \Gscale@div\@tempb{\linewidth}{\wd\pandoc@box}%
  \ifdim\@tempb\p@<\@tempa\p@\let\@tempa\@tempb\fi% select the smaller of both
  \ifdim\@tempa\p@<\p@\scalebox{\@tempa}{\usebox\pandoc@box}%
  \else\usebox{\pandoc@box}%
  \fi%
}
% Set default figure placement to htbp
\def\fps@figure{htbp}
\makeatother



\ifLuaTeX
\usepackage[bidi=basic,shorthands=off]{babel}
\else
\usepackage[bidi=default,shorthands=off]{babel}
\fi
\ifLuaTeX
  \usepackage{selnolig} % disable illegal ligatures
\fi


\setlength{\emergencystretch}{3em} % prevent overfull lines

\providecommand{\tightlist}{%
  \setlength{\itemsep}{0pt}\setlength{\parskip}{0pt}}



 


\KOMAoption{captions}{tableheading}
\makeatletter
\@ifpackageloaded{caption}{}{\usepackage{caption}}
\AtBeginDocument{%
\ifdefined\contentsname
  \renewcommand*\contentsname{Tabla de contenidos}
\else
  \newcommand\contentsname{Tabla de contenidos}
\fi
\ifdefined\listfigurename
  \renewcommand*\listfigurename{Listado de Figuras}
\else
  \newcommand\listfigurename{Listado de Figuras}
\fi
\ifdefined\listtablename
  \renewcommand*\listtablename{Listado de Tablas}
\else
  \newcommand\listtablename{Listado de Tablas}
\fi
\ifdefined\figurename
  \renewcommand*\figurename{Figura}
\else
  \newcommand\figurename{Figura}
\fi
\ifdefined\tablename
  \renewcommand*\tablename{Tabla}
\else
  \newcommand\tablename{Tabla}
\fi
}
\@ifpackageloaded{float}{}{\usepackage{float}}
\floatstyle{ruled}
\@ifundefined{c@chapter}{\newfloat{codelisting}{h}{lop}}{\newfloat{codelisting}{h}{lop}[chapter]}
\floatname{codelisting}{Listado}
\newcommand*\listoflistings{\listof{codelisting}{Listado de Listados}}
\captionsetup{labelsep=period}
\makeatother
\makeatletter
\makeatother
\makeatletter
\@ifpackageloaded{caption}{}{\usepackage{caption}}
\@ifpackageloaded{subcaption}{}{\usepackage{subcaption}}
\makeatother
\usepackage{bookmark}
\IfFileExists{xurl.sty}{\usepackage{xurl}}{} % add URL line breaks if available
\urlstyle{same}
\hypersetup{
  pdflang={es},
  colorlinks=true,
  linkcolor={black},
  filecolor={Maroon},
  citecolor={black},
  urlcolor={black},
  pdfcreator={LaTeX via pandoc}}


\author{}
\date{}
\begin{document}


\thispagestyle{empty}
\AddToHookNext{shipout/background}{
  \put(-4pt,-\dimexpr\paperheight+4pt\relax){
    \includegraphics[width=\dimexpr\paperwidth+8pt\relax,height=\dimexpr\paperheight+8pt\relax]{\detokenize{Imágenes y Recursos/Portada Informe Corporativo Blanco y Negro Minimalista.png}}
  }
}
\null
\newpage

\newgeometry{margin=2.5cm}

\tableofcontents
\newpage

\section{Resumen ejecutivo}\label{resumen-ejecutivo}

Completar en máximo una página.

Debe presentar una síntesis clara para una contraparte técnica o
institucional. Incluir brevemente el problema abordado, la fuente de
datos, el período analizado, el enfoque metodológico, los hallazgos
principales y las recomendaciones generales.

\section{Antecedentes, objetivo y
alcance}\label{antecedentes-objetivo-y-alcance}

\subsection{Antecedentes}\label{antecedentes}

Completar con el contexto de la asesoría, la necesidad de la contraparte
y la relevancia de analizar catálogos abiertos de terremotos.

\subsection{Objetivo general}\label{objetivo-general}

Completar con el objetivo principal del informe.

\subsection{Objetivos específicos}\label{objetivos-especuxedficos}

\begin{itemize}
\tightlist
\item
  Completar.
\item
  Completar.
\item
  Completar.
\end{itemize}

\subsection{Alcance del informe}\label{alcance-del-informe}

Completar indicando qué cubre el análisis y qué queda fuera. Recordar
que el informe no corresponde a predicción de terremotos, alerta sísmica
ni evaluación formal de amenaza.

\section{Descripción de datos y parámetros de
descarga}\label{descripciuxf3n-de-datos-y-paruxe1metros-de-descarga}

\subsection{Fuente de datos}\label{fuente-de-datos}

Completar con la fuente utilizada, formato de descarga y fecha de
obtención de los datos.

\subsection{Parámetros de descarga}\label{paruxe1metros-de-descarga}

Completar con período, magnitud mínima, zona geográfica, profundidad,
estado de revisión y formato del archivo.

\begin{longtable}[]{@{}lll@{}}
\caption{Parámetros de descarga utilizados para construir la base de
análisis.}\label{tbl-parametros-descarga}\tabularnewline
\toprule\noalign{}
Parámetro & Valor utilizado & Justificación \\
\midrule\noalign{}
\endfirsthead
\toprule\noalign{}
Parámetro & Valor utilizado & Justificación \\
\midrule\noalign{}
\endhead
\bottomrule\noalign{}
\endlastfoot
Fuente & completar & completar \\
Fecha de descarga & completar & completar \\
Período & completar & completar \\
Magnitud mínima & completar & completar \\
Zona geográfica & completar & completar \\
Formato & completar & completar \\
\end{longtable}

\subsection{Variables consideradas}\label{variables-consideradas}

Describir las variables disponibles y las variables derivadas que se
utilizarán en el análisis.

\begin{longtable}[]{@{}lll@{}}
\caption{Variables consideradas en el análisis
estadístico.}\label{tbl-variables}\tabularnewline
\toprule\noalign{}
Variable & Descripción & Uso en el análisis \\
\midrule\noalign{}
\endfirsthead
\toprule\noalign{}
Variable & Descripción & Uso en el análisis \\
\midrule\noalign{}
\endhead
\bottomrule\noalign{}
\endlastfoot
completar & completar & completar \\
completar & completar & completar \\
completar & completar & completar \\
\end{longtable}

\subsection{Depuración y control de
calidad}\label{depuraciuxf3n-y-control-de-calidad}

Completar con criterios de limpieza, revisión de duplicados, valores
faltantes, consistencia temporal, consistencia espacial y decisiones de
exclusión.

\section{Metodología y justificación
estadística}\label{metodologuxeda-y-justificaciuxf3n-estaduxedstica}

\subsection{Enfoque metodológico}\label{enfoque-metodoluxf3gico}

Completar con una explicación narrativa del enfoque utilizado. Evitar
describir solo comandos o salidas de software; explicar por qué cada
procedimiento es adecuado para responder el encargo.

\subsection{Indicadores descriptivos}\label{indicadores-descriptivos}

Completar con los indicadores que serán utilizados para caracterizar la
sismicidad.

\subsection{Visualizaciones}\label{visualizaciones}

Completar con las figuras previstas y su propósito.

\begin{longtable}[]{@{}lll@{}}
\caption{Visualizaciones previstas para comunicar los resultados del
informe.}\label{tbl-visualizaciones}\tabularnewline
\toprule\noalign{}
Figura & Propósito & Estado \\
\midrule\noalign{}
\endfirsthead
\toprule\noalign{}
Figura & Propósito & Estado \\
\midrule\noalign{}
\endhead
\bottomrule\noalign{}
\endlastfoot
completar & completar & pendiente \\
completar & completar & pendiente \\
completar & completar & pendiente \\
\end{longtable}

\subsection{Procedimientos
estadísticos}\label{procedimientos-estaduxedsticos}

Completar con pruebas, comparaciones, modelos o criterios inferenciales
si corresponde.

\subsection{Supuestos y criterios
metodológicos}\label{supuestos-y-criterios-metodoluxf3gicos}

Completar con los supuestos principales y su justificación.

\section{Resultados}\label{resultados}

\subsection{Caracterización general}\label{caracterizaciuxf3n-general}

Completar con la descripción general de la base analizada.

\subsection{Resultados temporales}\label{resultados-temporales}

Completar con tablas o figuras numeradas sobre evolución anual, mensual
o por período.

\subsection{Resultados espaciales}\label{resultados-espaciales}

Completar con mapas, distribución geográfica o comparación entre zonas.

\subsection{Magnitud y profundidad}\label{magnitud-y-profundidad}

Completar con resultados asociados a magnitud, profundidad y categorías
relevantes.

\subsection{Comparaciones principales}\label{comparaciones-principales}

Completar con las comparaciones solicitadas según el encargo.

\section{Discusión de resultados, supuestos, limitaciones y
sensibilidad}\label{discusiuxf3n-de-resultados-supuestos-limitaciones-y-sensibilidad}

\subsection{Interpretación de
resultados}\label{interpretaciuxf3n-de-resultados}

Completar conectando los resultados con la pregunta de asesoría.

\subsection{Supuestos y limitaciones}\label{supuestos-y-limitaciones}

Completar considerando completitud del catálogo, diferencias de
monitoreo, dependencia temporal por réplicas, tipo de magnitud,
incertidumbre de localización y sesgos por umbral.

\subsection{Análisis de sensibilidad}\label{anuxe1lisis-de-sensibilidad}

Completar con sensibilidad ante umbral de magnitud, período, definición
espacial, profundidad o estado de revisión.

\section{Conclusiones y
recomendaciones}\label{conclusiones-y-recomendaciones}

\subsection{Conclusiones}\label{conclusiones}

\begin{itemize}
\tightlist
\item
  Completar.
\item
  Completar.
\item
  Completar.
\end{itemize}

\subsection{Recomendaciones}\label{recomendaciones}

\begin{itemize}
\tightlist
\item
  Completar.
\item
  Completar.
\item
  Completar.
\end{itemize}

\section{Referencias}\label{referencias}

Completar con las citas utilizadas en el texto. Las referencias se
administran desde \texttt{referencias.bib}.

\section{Anexos reproducibles}\label{anexos-reproducibles}

\subsection{Parámetros de descarga}\label{paruxe1metros-de-descarga-1}

Completar con la URL, consulta, filtros o procedimiento utilizado para
obtener los datos.

\subsection{Código de procesamiento}\label{cuxf3digo-de-procesamiento}

Incluir código o enlaces a scripts solo como respaldo reproducible, no
como parte central de la narración del informe.

\subsection{Material complementario}\label{material-complementario}

Agregar tablas, figuras o detalles técnicos que respalden el informe sin
interrumpir la lectura principal.

\section{Bibliografía}\label{bibliografuxeda}




\end{document}
